\section{Estudio de Muros}

A continuación, se presenta el análisis detallado de los muros del edificio con el objetivo de establecer las características de su prediseño estructural.

\subsection{Largos de Muros}

Para determinar el espesor requerido en los elementos verticales, se midió la longitud total de los muros en las direcciones horizontal ($X$) y vertical ($Y$) para cada nivel. Los resultados obtenidos se resumen en la Tabla \ref{tab:largos_muros}:

\begin{table}[H]
    \centering
    \begin{tabular}{lcc}
        \toprule
        \textbf{Piso} & \textbf{Largo Horizontal [m]} & \textbf{Largo Vertical [m]} \\
        \midrule
        Subterráneo  & 81.22 & 63.76 \\
        Primer Piso  & 41.16 & 46.77 \\
        Segundo Piso & 36.41 & 38.83 \\
        Piso Tipo    & 37.31 & 45.20 \\
        \bottomrule
    \end{tabular}
    \caption{Longitudes totales de muros por nivel.}
    \label{tab:largos_muros}
\end{table}

\subsection{Corte Basal}

El esfuerzo de corte basal se determinó mediante la siguiente expresión:

\begin{equation}
    V = C \cdot I \cdot P
\end{equation}

Donde $C$ representa el coeficiente sísmico, $I$ el factor de importancia y $P$ el peso total de la estructura. Considerando los valores $C = 0.1$, $I = 1.0$ y $P = 9289.149 \text{ tonf}$, se obtiene un corte basal total de $V = 928.915 \text{ tonf}$.

\subsection{Índice de densidad de muros}

El índice de densidad de muros ($Im$) permite evaluar la rigidez lateral de la edificación en relación con su área en planta. La metodología de cálculo empleada es la siguiente:

\begin{equation}
    Im_i = \frac{L_i \cdot e_{prop}}{A_{planta}}
\end{equation}

Donde $L_i$ corresponde al largo acumulado de muros en la dirección analizada, $e_{prop}$ es el espesor propuesto y $A_{planta}$ es el área neta de la planta (descontando el área de la caja de ascensor de $7.47 \text{ m}^2$). Las áreas netas resultantes son: $1146.13 \text{ m}^2$ para el subterráneo, $319.93 \text{ m}^2$ para el primer piso, $322.53 \text{ m}^2$ para el segundo piso y $349.33 \text{ m}^2$ para el piso tipo.

Como referencia para el espesor de diseño, se utiliza la relación basada en el esfuerzo de corte admisible ($\tau$):

\begin{equation}
    e_{ref} = \frac{V}{\tau \cdot L_i}
\end{equation}

Se adopta un valor conservador de $\tau = 80 \text{ tonf/m}^2$. Los resultados del análisis se presentan a continuación:

\subsubsection{Índice de Densidad de Muros Horizontales}

\begin{table}[H]
    \centering
    \begin{tabular}{lccccc}
        \toprule
        \textbf{Nivel} & \textbf{Lx [m]} & \textbf{ex (ref) [m]} & \textbf{ex prop [m]} & \textbf{Área Planta [$m^2$]} & \textbf{Imx} \\
        \midrule
        Subterráneo  & 81.22 & 0.14 & 0.35 & 1146.13 & 2.5\% \\
        Primer Piso  & 41.16 & 0.28 & 0.25 & 319.93  & 3.2\% \\
        Segundo Piso & 36.41 & 0.32 & 0.35 & 322.53  & 4.0\% \\
        Piso Tipo    & 37.31 & 0.31 & 0.30 & 349.33  & 3.2\% \\
        \bottomrule
    \end{tabular}
    \caption{Resultados e índice de densidad para muros horizontales.}
    \label{tab:densidad_horiz}
\end{table}

\subsubsection{Índice de Densidad de Muros Verticales}

\begin{table}[H]
    \centering
    \begin{tabular}{lccccc}
        \toprule
        \textbf{Nivel} & \textbf{Ly [m]} & \textbf{ey (ref) [m]} & \textbf{ey prop [m]} & \textbf{Área Planta [$m^2$]} & \textbf{Imy} \\
        \midrule
        Subterráneo  & 141.62 & 0.08 & 0.35 & 1146.13 & 4.3\% \\
        Primer Piso  & 46.77  & 0.25 & 0.25 & 319.93  & 3.7\% \\
        Segundo Piso & 38.83  & 0.30 & 0.30 & 322.53  & 3.6\% \\
        Piso Tipo    & 45.20  & 0.26 & 0.20 & 349.33  & 2.6\% \\
        \bottomrule
    \end{tabular}
    \caption{Resultados e índice de densidad para muros verticales.}
    \label{tab:densidad_vert}
\end{table}

\subsection{Espesor Definitivo de Muros}

Dado que el índice de densidad ($Im$) se mantiene dentro del rango normativo/recomendado (entre el 2.5\% y 5.0\%) para ambas direcciones en todos los niveles, se validan los espesores propuestos. Por lo tanto, el diseño contempla un espesor de 0.35 m para el subterráneo, 0.25 m para el primer piso y 0.30 m para el segundo piso y piso tipo.